\chapter{Conclusion}\label{ch:conclusion}

This project has developed a Spatial Expected Threat (SxT) model that progresses through three increasingly sophisticated architectures, each addressing the key limitations of the previous version.

Version 1 established a working baseline using a Random Forest on positional and context features. While functional, it confirmed the known limitation of tabular models for this task: the model is dominated by a distance-to-goal heuristic, and even with match-state context features the absence of player positional data prevents it from capturing the richer spatial dynamics of a possession sequence.

Version 2 addressed this by moving to a dual-branch CNN architecture trained on StatsBomb 360 data, encoding the full pitch as a multi-channel spatial image. This allowed the model to see where all players are, not just the ball, and to reason about the spatial patterns of play. The resulting improvement in ROC AUC, from 0.68 to 0.83, represents the single largest gain in the project, though the CNN's implicit spatial encoding still limits its ability to reason explicitly about individual player relationships and formation shape.

Version 3 investigated whether adding a Transformer cross-attention layer on top of the frozen Version 2 CNN (encoding each player as an individual token and letting the Stage 1 threat logit query the player context) could improve on the CNN's aggregate spatial encoding. The results are positive and well-supported. A match-level bootstrap confirms that Version 3's ROC AUC gain of $+0.0066$ (0.830 $\to$ 0.836) is statistically significant (95\% CI $[+0.0014,\ +0.0125]$, $p = 0.017$). The initial calibration degradation (higher Brier Score and Log Loss for the raw Version 3 outputs) is fully correctable with post-hoc temperature scaling, after which Version 3 leads Version 2 on three of the four reported metrics. The scenario heatmaps show qualitatively that this gain comes from genuine formation-awareness: distinguishing a compact low block from a high press, responding to a 2v1 overload, and identifying an uncovered goal when the goalkeeper is off the line, in ways Version 2's density encoding cannot represent. An ablation further decomposes the gain, showing that a correction layer alone accounts for the majority of it, with explicit player-token attention contributing additional signal, a contribution that grows to 30\% of the total gain under the multi-query variant ($K=8$), which also achieves the highest AUC of any model in the project (0.8377). The architectural finding is therefore not a failure of the cross-attention design but a diagnosis of where the original single-query formulation is an information bottleneck, with a clear direction for improvement. The genuine, sizeable advance in the project remains the \emph{learned spatial threat surface} of Version 2; Version 3 builds on that foundation with a statistically reliable and architecturally interpretable improvement.

A controlled benchmark against VAEP~\cite{decroos2019}, trained and evaluated on the same StatsBomb split using the official \texttt{socceraction} implementation, places VAEP at ROC AUC 0.7359, above Version 1 and approximately 10 AUC points below Versions 2 and 3. This is the most direct available quantification, on a shared dataset, of the discriminative value of the 360 freeze-frame spatial encoding over the strongest tabular alternative.

The accompanying interactive frontend brings the model to life for a non-technical audience, allowing users to explore how player arrangements affect xT in real time.

In the taxonomy of non-shot threat models, our model occupies the gap between Singh's hand-crafted, player-blind zone grid and the dense-tracking EPV and pitch-control models~\cite{spearman2018,fernandez2019epv}: it learns a continuous threat surface (unlike Singh~\cite{singh2019}), conditions on player positions (unlike VAEP~\cite{decroos2019} and the event-only RL models~\cite{liu2020deep}), and does so from openly available freeze-frame snapshots rather than proprietary tracking data. The evaluation framework and controlled VAEP benchmark introduced here provide the first head-to-head comparison in this space on shared data, establishing a reproducible baseline against which future models using the same StatsBomb split can directly measure progress.

\section{Limitations and Future Work}

\begin{itemize}
    \item The 360 data is available for only a subset of matches and competitions, limiting the training set for Versions 2 and 3. Expanding this dataset through additional StatsBomb releases or alternative data providers would likely improve both versions.

    \item The class imbalance (goal rate $\approx 1.74\%$) makes training and evaluation challenging. More sophisticated imbalance handling, for example focal loss~\cite{lin2017}, could improve calibration further, particularly for Version 2. Focal loss modulates the per-example cross-entropy loss by a factor $(1-\hat{p})^\gamma$, which down-weights easy negatives more aggressively than the current \texttt{pos\_weight} capping strategy. The key hyperparameter to tune is $\gamma \geq 0$: at $\gamma = 0$ focal loss reduces to standard cross-entropy; values in the range $\gamma \in [0.5, 2.0]$ are typically explored first for heavily imbalanced binary tasks. For the 1.74\% positive rate here, larger $\gamma$ values suppress the gradient contribution of the many low-probability background events and could reduce the overconfidence that drives Version 2's relatively high Brier Score.

    \item Player identity and movement trajectories are not used; the model sees only a static snapshot at the moment of each event. Incorporating sequential information, for example via an LSTM or temporal Transformer, could capture the momentum and direction of play.

    \item The end-position inference rules used during visualisation and frontend evaluation remain hand-crafted. Learning these from data, for example by training a model to predict the empirical end position given the ball location and player context, would remove the last fixed constants from the evaluation pipeline and potentially produce more realistic per-action threat surfaces.

    \item The architectural experiments in Section~\ref{sec:v3_arch_experiments} demonstrate that multi-query cross-attention ($K=8$) modestly but consistently improves on the single-query design (AUC $0.8363 \to 0.8377$), confirming that the original single scalar-logit query is an information bottleneck. A natural next step is to investigate whether larger $K$, query conditioning on both the Stage 1 logit and a richer ball-context vector, or a full cross-attention stack (rather than a single layer) can unlock further player-context signal , subject to the regularisation challenges that already cause calibration to degrade as Stage 2 becomes more expressive.
\end{itemize}

The framework is designed to be extensible: additional data sources, richer player encodings, or more powerful Transformer architectures can be slotted into the same evaluation pipeline without requiring changes to the comparison methodology.
